\section{Evaluation}
\label{sec:evaluation}

We report the case's metrics: coverage (share of the total heat demand covered by server heat; the
rest comes from the backup heater), dumped heat, utilisation (server running hours $\div$ (servers
$\times$ 8{,}760)), missed deadlines in per cent of all jobs (fleet runs) or of Type~1 jobs (network experiments), and the work moved between sites.
Commands and source files are listed in Appendix~\ref{app:repro}; network CSV rows also
record their exact command lines.

\subsection{Reference house}
\label{sec:eval-house}

Table~\ref{tab:house-fleet} separates what the house could absorb from what it absorbs under
sourced demand. With unlimited work, four servers cover \qty{58.05}{\percent} of the heat demand at
\qty{64.7}{\percent} utilisation without dumping. This is the physical limit of the house; the
fleet model reproduces it within \qty{0.02}{\percent} (\num{22657.0} against \num{22661.0}
server-hours) when it is given a backlog that never runs out. With the sourced job demand, scaled
to one house, the same four servers run \num{8917} server-hours: coverage falls to
\qty{22.8}{\percent} and utilisation to \qty{25.4}{\percent}, still without dumping. Of
\num{77840} jobs, 6,905 (\qty{8.9}{\percent}) miss their deadline, and \num{5763} of the
\num{5771} jobs that never finished arrived between May and September: with only hot water to heat, the
scheduler has no useful heat to run them on. Four servers is the case's comparison count, not an
optimised one (Section~\ref{sec:eval-sizing}).

\begin{table}[ht]
  \centering
  \caption{Reference house and real house in the fleet model (Basel 2025, L4 hardware profile). The
    fleet columns use the same jobs and seed; the unlimited-work column is the case's physical
    limit from the year model (cross-checked by the fleet model).}
  \label{tab:house-fleet}
  \small
  \setlength{\tabcolsep}{4pt}
  \begin{tabular}{lrrrr}
    \toprule
    & \multicolumn{3}{c}{Reference house} & Real house \\
    \cmidrule(lr){2-4}\cmidrule(lr){5-5}
    & unlimited work & 4 servers & 1 server & 4 servers \\
    \midrule
    Heat demand [kWh] & \num{31231} & \num{31231} & \num{31231} & \num{41356} \\
    Heat demand covered & \qty{58.05}{\percent} & \qty{22.8}{\percent} & \qty{18.9}{\percent} & \qty{21.6}{\percent} \\
    Heat dumped [kWh] & 0 & 0 & 0 & 0 \\
    Utilisation & \qty{64.7}{\percent} & \qty{25.4}{\percent} & \qty{84.1}{\percent} & \qty{31.9}{\percent} \\
    Server-hours sold & \num{22661} & \num{8917} & \num{7366} & \num{11164} \\
    Missed deadlines & -- & \qty{8.9}{\percent} & \qty{5.9}{\percent} & \qty{0.01}{\percent} \\
    Computing income [CHF] & -- & \num{13360} & \num{13333} & \num{13780} \\
    Net after hardware [CHF/year] & -- & \num{-14317} & \num{5252} & \num{-14406} \\
    Case break-even [CHF/server-h] & 0.2261 & 0.2261 & 0.2261 & 0.2261 \\
    \bottomrule
  \end{tabular}
\end{table}

\paragraph{Real house.} The organisers provided the monthly gas readings of a real, anonymised
house that heats with gas. We convert January to September 2025 (later months are missing) to heat
with the 86\% minimum seasonal efficiency of new gas boilers \cite{eu813}, and fit the case
formula's two house parameters by least squares on those nine months. The fit gives
\qty{10.13}{kW} design heating power, against the case's \qty{10}{kW}, and \qty{38.8}{kWh} of heat
per day outside the heating season, against the case's \qty{12}{kWh} of hot water ($R^2 = 0.996$,
monthly errors $-5.9$ to $+8.4$\%). The case formula therefore describes this house's space heating
well; its summer use, which includes circulation and boiler losses that monthly data cannot
separate, is three times the case's hot water. Run with the same jobs, the real house misses 11
deadlines instead of 6,905 (Table~\ref{tab:house-fleet}): its summer heat absorbs the summer work.
Its income rises only slightly, because the extra work is mostly low-priced Type~2. The boiler's
real efficiency and age are unknown; the 86\% is an upper value for a new appliance.

\paragraph{Planning on one house.} The following comparison uses the network model with unlimited
work (commit \texttt{1863f57}). On a single house the greedy rule B1 is optimal: it reaches the full-year upper bound exactly, and
the rolling-horizon LP cannot improve on it (Table~\ref{tab:house-strategies}). The reason is that on
the reference house the minimum load is zero, the tank is lossless and work is unlimited, so using
heat now never reduces what can be used later; four servers are the limit in cold hours whatever the
strategy. Always-on (B0) reaches the same coverage but only adds \qty{9903}{kWh} of dumped heat in
mild hours. Planning on imperfect forecasts costs little: coverage falls from \qty{58.05}{\percent}
to \qty{57.82}{\percent} at $\sigma_{24} = \qty{20}{\percent}$, because the LP plans the current hour
for the forecast world and its plan can be slightly worse in reality.
\emph{On one house, intelligence does not pay; heat physics decides.}

\begin{table}[ht]
  \centering
  \caption{Reference-house strategies (four servers, \qty{35}{kWh} tank).
    Forecast errors are synthetic AR(1) innovations ($\rho=0.8$); UB is a bound.}
  \label{tab:house-strategies}
  \small
  \setlength{\tabcolsep}{4pt}
  \begin{tabular}{llrrrr}
    \toprule
    Strategy & Information & Coverage & Dumped [kWh] & Server-hours & Utilisation \\
    \midrule
    B0 always on & current hour & \qty{58.05}{\percent} & \num{9903} & \num{35040} & \qty{100}{\percent} \\
    B1 local greedy & current hour & \qty{58.05}{\percent} & 0 & \num{22661.0} & \qty{64.67}{\percent} \\
    LP rolling horizon & perfect foresight & \qty{58.05}{\percent} & 0 & \num{22661.0} & \qty{64.67}{\percent} \\
    LP rolling horizon & error $\sigma_{24} = \qty{5}{\percent}$ & \qty{58.03}{\percent} & 0 & \num{22654.7} & \qty{64.65}{\percent} \\
    LP rolling horizon & error $\sigma_{24} = \qty{10}{\percent}$ & \qty{57.99}{\percent} & 0 & \num{22639.1} & \qty{64.61}{\percent} \\
    LP rolling horizon & error $\sigma_{24} = \qty{20}{\percent}$ & \qty{57.82}{\percent} & 0 & \num{22572.9} & \qty{64.42}{\percent} \\
    UB full-year LP & perfect foresight, year & \qty{58.05}{\percent} & 0 & \num{22661.0} & \qty{64.67}{\percent} \\
    \bottomrule
  \end{tabular}
\end{table}

\subsection{Sizing}
\label{sec:eval-sizing}

Each additional server adds fewer heat-feasible hours. With unlimited work, the first server can run
\num{7631} hours a year, the fourth \num{4882}, the sixth \num{3796} and the thirteenth only 5.
Under sourced demand the binding limit is different: at the scaled demand, one to four servers
earn almost the same income (CHF~13,300 to 14,100), because the work is limited by demand in winter
and by heat in summer. One server then has the only positive net after hardware
(Table~\ref{tab:breakeven}): $+$CHF~5,252 a year, against $-$CHF~14,317 for four. Its coverage is
lower (\qty{18.9}{\percent} instead of \qty{22.8}{\percent}). It also misses fewer deadlines
(\qty{5.9}{\percent}) while rejecting more Type~2 work on price (1,250 jobs instead of 878): with
less capacity the load-based price quotes higher, so more summer heat remains for Type~1 work. Four servers would break even at 2.8 times the scaled demand. These results depend on the
demand baseline, a labelled scenario value (Section~\ref{sec:limits}).

\subsection{Local scenario: a pool as a summer heat sink}
\label{sec:eval-local}

A local fleet scenario places five case houses (four servers each, at the Münchenstein electricity
price of CHF~0.3198/kWh) around FHNW Campus Dreispitz and one indoor pool (20 servers) with the
annual heat demand of the renovated Allschwil indoor pool, \qty{627000}{kWh}, spread flat over the
year. It uses sourced, unscaled Basel job volumes, and customers accept every Type~2 price
(elasticity 0). Positions, the flat profile and the pool's 20 servers are scenario values. The same heat-first policy places work where the free share of the hourly
heat budget is largest. In January the pool does 34\% of the work; in July, when the houses only
need hot water, it does 76\%. Over the year 125 of \num{778536} jobs (\qty{0.02}{\percent}) miss
their deadline. The houses cover about 30\% of their heat, the pool 7.9\%. Every site's net after
hardware is still negative at four L4 servers per house (about $-$CHF~12,600 per house).

\subsection{Strategy comparison on a Swiss test network}
\label{sec:eval-network}

The synthetic experiments in this and the next subsections were computed with the research
strategies of Section~\ref{sec:scheduling} at commit \texttt{1863f57}. The test network has 33 sites
in nine Swiss cities:
the reference house, single- and multi-family houses, nine hotels with spa and five indoor pools.
Locations and sizes are \emph{synthetic test data}; heat demand comes from the case formula
(houses) or from the profiles named in the scenario file. All results in this subsection are
\emph{scenarios}, not forecasts of a real network.

At $\varphi=1$, B1 and B2 each cover \qty{83.42}{\percent} of demand with no dumped heat;
B2-reserve covers \qty{83.40}{\percent}. With scarce filler, placement also changes coverage
(Table~\ref{tab:exp-b}). The strategies
differ in missed deadlines (Table~\ref{tab:exp-a}): with \qty{30}{\percent} paid work, B1 misses
\qty{15.69}{\percent} of jobs, while B2 misses none by moving about \num{103000} server-hours to
sites with summer heat demand.

\begin{table}[ht]
  \centering
  \caption{Experiment A: paid-work share at 33 synthetic sites, $\varphi=1$.
    Greedy coverage is \qty{83.42}{\percent} with zero dumped heat
    (B2-reserve: \qty{83.40}{\percent}). The work share is a scenario input.}
  \label{tab:exp-a}
  \small
  \begin{tabular}{rrrrrrr}
    \toprule
    Type 1 & B1 & B2 & B2 moved & B2-reserve & LP & LP moved \\
    share & missed & missed & [server-h] & moved [server-h] & missed & [server-h] \\
    \midrule
    \qty{10}{\percent} & \qty{0}{\percent} & \qty{0}{\percent} & \num{7725} & \num{7725} & \qty{0}{\percent} & \num{0} \\
    \qty{30}{\percent} & \qty{15.69}{\percent} & \qty{0}{\percent} & \num{102754} & \num{64230} & \qty{0}{\percent} & \num{49884} \\
    \qty{50}{\percent} & \qty{25.19}{\percent} & \qty{0}{\percent} & \num{205267} & \num{161968} & \qty{0}{\percent} & \num{149009} \\
    \qty{70}{\percent} & \qty{29.47}{\percent} & \qty{0}{\percent} & \num{309325} & \num{265260} & \qty{0}{\percent} & \num{253300} \\
    \bottomrule
  \end{tabular}
\end{table}

\subsection{Why the network helps}
\label{sec:eval-why}

Figures~\ref{fig:f10} and \ref{fig:f11} show the deadline and routing effects at a glance.
Their underlying scenario values are reported in Tables~\ref{tab:exp-a} and \ref{tab:exp-d}.
% Data from figures/data/a_type1_share.dat, produced by
% python3 scripts/paper_figure_data.py from data/derived/network/experiments/A_type1_share[-lp].csv
\begin{figure}[ht]
  \centering
  \begin{tikzpicture}
    \begin{axis}[name=left, width=0.47\textwidth, height=5.2cm, xlabel={Type~1 share [\%]},
        ylabel={Missed Type~1 deadlines [\%]}, xtick={10,30,50,70}, ymin=0, ymax=35,
        legend style={font=\scriptsize, at={(0.02,0.98)}, anchor=north west}, unbounded coords=jump]
      \addplot+[mark=square*] table[x=share, y=B1missed] {figures/data/a_type1_share.dat};
      \addplot+[mark=triangle*] table[x=share, y=B2missed] {figures/data/a_type1_share.dat};
      \addplot+[mark=*, dashed] table[x=share, y=LPmissed] {figures/data/a_type1_share.dat};
      \legend{B1 local, B2 network, LP}
    \end{axis}
    \begin{axis}[at={(left.east)}, anchor=west, xshift=1.6cm, width=0.47\textwidth, height=5.2cm,
        xlabel={Type~1 share [\%]}, ylabel={Moved Type~1 work [$10^3$ server-h]}, xtick={10,30,50,70},
        ymin=0, legend style={font=\scriptsize, at={(0.02,0.98)}, anchor=north west}, unbounded coords=jump]
      \addplot+[mark=triangle*] table[x=share, y=B2moved] {figures/data/a_type1_share.dat};
      \addplot+[mark=diamond*] table[x=share, y=B2Rmoved] {figures/data/a_type1_share.dat};
      \addplot+[mark=*, dashed] table[x=share, y=LPmoved] {figures/data/a_type1_share.dat};
      \legend{B2, B2-reserve, LP}
    \end{axis}
  \end{tikzpicture}
  \caption{Experiment A at 33 synthetic sites ($\varphi=1$): missed Type~1 jobs (left)
    and moved server-hours (right). Greedy heat coverage is \qty{83.42}{\percent}
    except B2-reserve (\qty{83.40}{\percent}); no heat is dumped.}
  \label{fig:f10}
\end{figure}

% Figure F11 (T-12): data from figures/data/d_network_mix.dat (python3 scripts/paper_figure_data.py)
\begin{figure}[ht]
  \centering
  \begin{tikzpicture}
    \begin{axis}[ybar, bar width=12pt, width=0.6\textwidth, height=5cm, ymin=0, ymax=20,
        ylabel={Missed Type~1 deadlines [\%]}, symbolic x coords={all,residential}, xtick=data,
        xticklabels={{all 33 sites}, {residential only (19)}}, enlarge x limits=0.4,
        legend style={font=\scriptsize, at={(0.02,0.98)}, anchor=north west}, nodes near coords,
        nodes near coords style={font=\scriptsize}, every node near coord/.append style={/pgf/number format/fixed, /pgf/number format/precision=1}]
      \addplot table[x=network, y=B1missed] {figures/data/d_network_mix.dat};
      \addplot table[x=network, y=B2missed] {figures/data/d_network_mix.dat};
      \addplot table[x=network, y=LPmissed] {figures/data/d_network_mix.dat};
      \legend{B1 local, B2 network, LP}
    \end{axis}
  \end{tikzpicture}
  \caption{Experiment D: missed Type~1 jobs at \qty{30}{\percent} paid-work share,
    $\varphi=1$. The synthetic residential-only network has no year-round heat sinks.}
  \label{fig:f11}
\end{figure}


In summer a house can only absorb heat for hot water, about 0.6 servers per day in the case. Paid
work above that cannot run at home without dumping heat. If the network contains year-round heat
sinks (hotels with spas, indoor pools), the work moves there and no deadline is missed. Remove them
and moving still happens, but \qty{13.80}{\percent} of jobs are missed, because every house is short
at the same time (Table~\ref{tab:exp-d}).

\begin{table}[ht]
  \centering
  \caption{Experiment D: network mix at \qty{30}{\percent} Type~1 share and $\varphi=1$.
    All sites are synthetic; B1 and B2 have the same heat coverage.}
  \label{tab:exp-d}
  \small
  \begin{tabular}{lrrrrrrr}
    \toprule
    Network & Sites & B1 & B2 & B2 moved & Coverage & LP & LP moved \\
    & & missed & missed & [server-h] & & missed & [server-h] \\
    \midrule
    all sites & 33 & \qty{15.69}{\percent} & \textbf{\qty{0}{\percent}} & \num{102754} & \qty{83.42}{\percent} & \qty{0}{\percent} & \num{49884} \\
    residential only & 19 & \qty{15.69}{\percent} & \textbf{\qty{13.80}{\percent}} & \num{63333} & \qty{63.18}{\percent} & \qty{13.07}{\percent} & \num{3001} \\
    \bottomrule
  \end{tabular}
\end{table}

Timing at home does not substitute for the network (Table~\ref{tab:reserve}). Keeping tank room for
tomorrow's work (B1-reserve) saves only 0.07 points of missed deadlines (\qty{15.62}{\percent}
versus \qty{15.69}{\percent} for B1), because a daily shortfall returns the next day. Combined with moves (B2-reserve) it reduces the work that has to move by
\qty{37}{\percent}, at a cost of 0.02 points of coverage, because origins keep tank room and need to
send less work away.

\begin{table}[ht]
  \centering
  \caption{Local timing versus moving work.
    Type~1 share \qty{30}{\percent}, $\varphi=1$, 33 synthetic sites.}
  \label{tab:reserve}
  \begin{tabular}{lrrrr}
    \toprule
    Strategy & Coverage & Dumped [kWh] & Missed Type 1 & Moved Type 1 [server-h] \\
    \midrule
    B1 & \qty{83.42}{\percent} & 0 & \qty{15.69}{\percent} & 0 \\
    B1-reserve & \qty{83.42}{\percent} & 0 & \qty{15.62}{\percent} & 0 \\
    B2 & \qty{83.42}{\percent} & 0 & \qty{0}{\percent} & \num{102754} \\
    B2-reserve & \qty{83.40}{\percent} & 0 & \qty{0}{\percent} & \num{64230} \\
    \bottomrule
  \end{tabular}
\end{table}

\paragraph{Transfer delay.} The results above move work without friction. Experiment~E
lets moved work start at the destination only $\delta$ hours after the cohort's
arrival (``data staged at arrival'', a lower bound on friction). The delays are an assumption
sweep; no sourced delay exists. Even $\delta = \qty{12}{h}$ costs no deadline: B2 stays at
\qty{0}{\percent} missed (Table~\ref{tab:exp-e}), because the year-round sinks have enough free
summer capacity to finish the excess in the remaining \qty{12}{h}. B2 even moves less with a delay,
from \num{102754} to \num{76680} and \num{61727} server-hours, i.e.\ from 1.57 to 0.94 moved
server-hours per avoided missed job: without a delay, its persistence test sends work away that would
still have finished at home. With $\delta \ge \qty{24}{h}$ (the deadline) no work can move and missed
deadlines return to B1's level.

\begin{table}[ht]
  \centering
  \caption{Experiment E: assumed transfer delay $\delta$ (33 synthetic sites).
    ``Avoided jobs'' compares each strategy with B1 (\num{65352} missed jobs).
    The LP at $\delta=0$ appears in Table~\ref{tab:exp-a}; $\delta=1$~h was not run with the LP.}
  \label{tab:exp-e}
  \small
  \setlength{\tabcolsep}{4pt}
  \begin{tabular}{rrrrrrr}
    \toprule
    $\delta$ & B2 & B2 moved & B2 moved per & B2-reserve & B2-reserve moved & LP missed \\
    & missed & [server-h] & avoided job & moved [server-h] & per avoided job & / moved \\
    \midrule
    \qty{0}{h} & \qty{0}{\percent} & \num{102754} & 1.57 & \num{64230} & 0.98 & Table~\ref{tab:exp-a} \\
    \qty{1}{h} & \qty{0}{\percent} & \num{102771} & 1.57 & \num{64230} & 0.98 & -- \\
    \qty{4}{h} & \qty{0}{\percent} & \num{76680} & 1.17 & \num{60527} & 0.93 & \qty{0}{\percent} / \num{49793} \\
    \qty{12}{h} & \qty{0}{\percent} & \num{61727} & 0.94 & \num{58775} & 0.90 & \qty{0}{\percent} / \num{49847} \\
    \bottomrule
  \end{tabular}
\end{table}

\subsection{Sensitivity}
\label{sec:eval-sensitivity}

\paragraph{Paid-work volume.} From a Type~1 share of \qty{10}{\percent} to \qty{70}{\percent}, B1's
missed deadlines grow from \qty{0}{\percent} to \qty{29.47}{\percent}; B2 stays at
\qty{0}{\percent}, while the moved server-hours grow roughly linearly (Table~\ref{tab:exp-a}).

\paragraph{Filler availability.} When the filler availability $\varphi$ falls from 1.0 to 0.2,
coverage falls from \qty{83.42}{\percent} to about \qty{21}{\percent} for every strategy
(Table~\ref{tab:exp-b}). A shortage of sellable work, not the algorithm, limits the delivered heat.
B2 helps coverage a little when filler is scarce (0.5 to 0.9 points), because moved Type~1 work
runs where heat is needed and filler would otherwise be missing.

\begin{table}[ht]
  \centering
  \caption{Experiment B: filler availability $\varphi$ at \qty{30}{\percent} Type~1 share
    (33 synthetic sites). The $\varphi=0.35$ scenario is the filler-capacity analogy of the demand
    research in the repository (\texttt{docs/research/demand.md}).}
  \label{tab:exp-b}
  \begin{tabular}{rrrrr}
    \toprule
    $\varphi$ & B1 coverage & B2 coverage & B1 missed & B2 missed \\
    \midrule
    1.0 & \qty{83.42}{\percent} & \qty{83.42}{\percent} & \qty{15.69}{\percent} & \qty{0}{\percent} \\
    0.6 & \qty{54.21}{\percent} & \qty{54.71}{\percent} & \qty{15.54}{\percent} & \qty{0}{\percent} \\
    0.35 & \qty{34.58}{\percent} & \qty{35.38}{\percent} & \qty{14.93}{\percent} & \qty{0}{\percent} \\
    0.2 & \qty{21.41}{\percent} & \qty{22.34}{\percent} & \qty{14.82}{\percent} & \qty{0}{\percent} \\
    \bottomrule
  \end{tabular}
\end{table}

\paragraph{Forecast error in the network.} In the synthetic error sweep, LP coverage falls
from \qty{83.42}{\percent} to \qty{83.23}{\percent} at
$\sigma_{24}=\qty{20}{\percent}$, with no missed deadlines or dumped heat
(Table~\ref{tab:exp-c}). This tests the stated error process rather than a measured forecast.

\begin{table}[ht]
  \centering
  \caption{Experiment C: LP forecast-error sensitivity (33 synthetic sites,
    \qty{30}{\percent} Type~1 share, $\varphi=1$). AR(1) parameters $\rho=0.8$ and
    $\sigma_{24}$ are uncalibrated; the perfect-foresight row is from experiment~A.}
  \label{tab:exp-c}
  \begin{tabular}{lrrrr}
    \toprule
    Forecast & Coverage & Dumped [kWh] & Missed Type 1 & Moved [server-h] \\
    \midrule
    perfect foresight & \qty{83.42}{\percent} & 0 & \qty{0}{\percent} & \num{49884} \\
    $\sigma_{24} = 0.05$ & \qty{83.39}{\percent} & 0 & \qty{0}{\percent} & \num{49882} \\
    $\sigma_{24} = 0.10$ & \qty{83.35}{\percent} & 0 & \qty{0}{\percent} & \num{49929} \\
    $\sigma_{24} = 0.20$ & \qty{83.23}{\percent} & 0 & \qty{0}{\percent} & \num{49967} \\
    \bottomrule
  \end{tabular}
\end{table}

\paragraph{Economic sensitivities.} Compute price, hardware, electricity price and heat fee: see
Section~\ref{sec:economics}.

\subsection{Summary of findings}

\begin{enumerate}
  \item[(i)] With unlimited work four servers cover 58\% of the reference house's heat; with sourced
    demand they cover 22.8\%, dump nothing and miss 8.9\% of deadlines, mostly in summer.
  \item[(ii)] At the sourced demand, hardware decides the result: one server covers its cost,
    four need 2.8 times the demand or 2.6 times the Type~1 price.
  \item[(iii)] Assuming 86\% boiler efficiency, the case formula fits a real house's space heating; the real house's larger summer
    heat use removes most of the summer deadline misses.
  \item[(iv)] On one house the greedy rule reaches the full-year upper bound. Value from routing
    arises in a network with year-round heat sinks: the executable policy moves 76\% of July work
    to the Dreispitz pool, and in the synthetic experiments greedy routing (B2) already removes the
    missed deadlines; the LP mainly reduces the work moved.
  \item[(v)] The limits of these statements are listed in Section~\ref{sec:limits}.
\end{enumerate}
